\subsection{Hasil Kali Titik}

Terdapat dua cara mengalikan vektor yang sangat penting dalam
penerapan. Cara pertama disebut \textbf{hasil kali titik}. Ketika kita mengambil hasil kali titik vektor, hasilnya
adalah sebuah skalar. Karena alasan ini, hasil kali titik juga
disebut \textbf{hasil kali skalar} dan terkadang \textbf{hasil kali dalam}. Definisinya adalah sebagai berikut.
\index{hasil kali titik}
\index{hasil kali skalar}
\index{hasil kali dalam}

\begin{definition}{Hasil Kali Titik}\label{def:dotproductdefn}
Misalkan $\vect{u},\vect{v}$ adalah dua vektor dalam $\mathbb{R}^{n}$. Maka kita
mendefinisikan \textbf{hasil kali titik} $\vect{u}\dotprod \vect{v}$ sebagai
\begin{equation*}
\vect{u}\dotprod \vect{v} = \sum_{k=1}^{n}u_{k}v_{k}
\end{equation*}
\end{definition}

Hasil kali titik $\vect{u}\dotprod \vect{v}$ terkadang dinyatakan sebagai $(\vect{u},\vect{v})$, dengan koma menggantikan $\dotprod $. Hasil kali ini juga dapat ditulis
sebagai $\left\langle \vect{u},\vect{v}\right\rangle $. Jika kita menulis vektor sebagai matriks kolom atau baris, hasil kali tersebut
sama dengan hasil kali matriks $\vect{v}\vect{w}^{T}$.

Perhatikan contoh berikut.

\begin{example}{Menghitung Hasil Kali Titik}\label{exa:dotproduct}
Tentukan $\vect{u} \dotprod \vect{v}$ untuk
\begin{equation*}
\vect{u}
=
\begin{bmatrix}
1 \\
2 \\
0 \\
-1
\end{bmatrix},
\vect{v}
=
\begin{bmatrix}
0 \\
1 \\
2 \\
3
\end{bmatrix}
\end{equation*}
\end{example}

\begin{solution}
Berdasarkan Definisi \ref{def:dotproductdefn}, kita harus menghitung
\begin{equation*}
\vect{u}\dotprod \vect{v} = \sum_{k=1}^{4}u_{k}v_{k}
\end{equation*}

Hasilnya diberikan oleh
\begin{eqnarray*}
\vect{u} \dotprod \vect{v}
&=&
(1)(0) + (2)(1) + (0)(2) + (-1)(3) \\
&=&
0 + 2 + 0 + -3 \\
&=&
-1
\end{eqnarray*}
\end{solution}

Berdasarkan definisi ini, terdapat beberapa sifat penting yang dipenuhi
oleh hasil kali titik.
\index{hasil kali titik!sifat-sifat}

\begin{proposition}{Sifat-Sifat Hasil Kali Titik}\label{prop:propertiesdotproduct}
Misalkan $k $ dan $p $ menyatakan skalar, serta $\vect{u},\vect{v},\vect{w}$ menyatakan vektor.
Maka hasil kali titik $\vect{u} \dotprod \vect{v}$ memenuhi sifat-sifat berikut.
\begin{itemize}
\item
$\vect{u}\dotprod \vect{v}= \vect{v}\dotprod \vect{u} $
\item
$\vect{u}\dotprod \vect{u}\geq 0 \text{ dan sama dengan nol jika dan hanya jika }\vect{u}=\vect{0}$
\item
$\left( k\vect{u}+p\vect{v}\right) \dotprod \vect{w}=k\left( \vect{u}\dotprod \vect{w}\right) +p\left( \vect{v}\dotprod \vect{w}\right)$
\item
$\vect{u}\dotprod\left( k\vect{v}+p\vect{w}\right) =k\left(
\vect{u}\dotprod \vect{v}\right) +p\left( \vect{u}\dotprod \vect{w}\right)$
\item
$\vectlength \vect{u}\vectlength ^{2}=\vect{u}\dotprod \vect{u} $
\end{itemize}
\end{proposition}

Buktinya diserahkan sebagai latihan. Proposisi ini memberi tahu kita bahwa kita juga dapat menggunakan hasil kali titik untuk mencari panjang suatu vektor.

\begin{example}{Panjang Vektor}\label{exa:dotproductlength}
Tentukan panjang
\begin{equation*}
\vect{u}
=
\leftB
\begin{array}{r}
2 \\
1 \\
4 \\
2
\end{array}
\rightB
\end{equation*}
Dengan kata lain, tentukan $\vectlength \vect{u} \vectlength .$ 
\end{example}

\begin{solution}
Berdasarkan Proposisi \ref{prop:propertiesdotproduct}, $\vectlength \vect{u} \vectlength ^{2} = \vect{u} \dotprod \vect{u}$.
Oleh karena itu, $\vectlength \vect{u} \vectlength = \sqrt {\vect{u} \dotprod \vect{u}}$.
Pertama, hitung $\vect{u} \dotprod \vect{u}$.

Hasilnya diberikan oleh
\begin{eqnarray*}
\vect{u} \dotprod \vect{u}
&=&
(2)(2) + (1)(1) + (4)(4) + (2)(2) \\
&=&
4 + 1 + 16 + 4 \\
&=&
25
\end{eqnarray*}

Kemudian,
\begin{eqnarray*}
\vectlength \vect{u} \vectlength
&=& \sqrt {\vect{u} \dotprod \vect{u}} \\
&=& \sqrt{25} \\
&=& 5
\end{eqnarray*}
\end{solution}

Anda mungkin ingin membandingkannya dengan definisi panjang kita sebelumnya,
yang diberikan dalam Definisi \ref{def:lengthofvector}.

\textbf{Ketaksamaan Cauchy--Schwarz} adalah ketaksamaan mendasar yang dipenuhi oleh hasil kali titik.
\index{ketaksamaan Cauchy--Schwarz} Ketaksamaan ini diberikan dalam teorema berikut.

\begin{theorem}{Ketaksamaan Cauchy--Schwarz}\label{thm:cauchyschwarzinequality}
Hasil kali titik memenuhi ketaksamaan
\begin{equation}
\left\vert \vect{u}\dotprod \vect{v}\right\vert \leq \vectlength \vect{u}\vectlength \vectlength \vect{v}\vectlength   \label{cauchy}
\end{equation}
Lebih lanjut, kesamaan diperoleh jika dan hanya jika salah satu dari $\vect{u}$ atau $\vect{v}$ merupakan kelipatan skalar dari yang lain.
\end{theorem}

\ifdefined\showproofs
\begin{proof}
Pertama, perhatikan bahwa jika $\vect{v}=\vect{0}$, kedua ruas \ref{cauchy}
sama dengan nol sehingga ketaksamaan berlaku dalam kasus ini. Oleh karena itu, dalam pembahasan berikut
kita akan mengasumsikan bahwa $\vect{v}\neq \vect{0}$.

Definisikan sebuah fungsi dari $t\in \mathbb{R}$ dengan
\begin{equation*}
f\left( t\right) =\left( \vect{u}+t\vect{v}\right) \dotprod \left( \vect{u}+
t\vect{v}\right)
\end{equation*}
Maka, berdasarkan Proposisi \ref{prop:propertiesdotproduct}, $f\left( t\right) \geq 0$ untuk semua $t\in \mathbb{R}$.
Juga dari Proposisi \ref{prop:propertiesdotproduct},
\begin{eqnarray*}
f\left( t\right) &=&\vect{u}\dotprod \left( \vect{u}+t\vect{v}\right) +
t\vect{v}\dotprod \left( \vect{u}+t\vect{v}\right) \\
&=&\vect{u}\dotprod \vect{u}+t\left( \vect{u}\dotprod \vect{v}\right) + t \vect{v}\dotprod \vect{u}+
t^{2}\vect{v}\dotprod \vect{v} \\
&=&\vectlength \vect{u}\vectlength ^{2}+2t\left( \vect{u}\dotprod \vect{v}\right) +\vectlength
\vect{v}\vectlength ^{2}t^{2}
\end{eqnarray*}

Ini berarti grafik $y=f\left( t\right) $ adalah parabola yang membuka
ke atas dan titik puncaknya menyentuh sumbu $t$, atau seluruh grafik berada
di atas sumbu $t$. Dalam kasus pertama, terdapat suatu $t$ sehingga $f\left(
t\right) =0$, dan hal ini mengharuskan $\vect{u}+t\vect{v}=\vect{0}$ sehingga satu vektor merupakan
kelipatan vektor lainnya. Maka jelas bahwa kesamaan berlaku dalam \ref{cauchy}. Dalam
kasus ketika $\vect{v}$ bukan kelipatan $\vect{u}$, diperoleh
$f\left( t\right) >0$ untuk semua $t$, yang berarti $f\left( t\right) $ tidak memiliki akar
real sehingga, dari rumus kuadrat,
\begin{equation*}
\left( 2\left( \vect{u}\dotprod \vect{v}\right) \right) ^{2}-4\vectlength \vect{u}
\vectlength ^{2}\vectlength \vect{v}\vectlength ^{2}<0
\end{equation*}
yang setara dengan $\left\vert \vect{u}\dotprod \vect{v} \right\vert
<\vectlength \vect{u}\vectlength \vectlength \vect{v}\vectlength $.
\end{proof}
\fi
Perhatikan bahwa bukti ini hanya didasarkan pada sifat-sifat
hasil kali titik yang tercantum dalam Proposisi \ref{prop:propertiesdotproduct}. Ini berarti bahwa
setiap kali suatu operasi memenuhi sifat-sifat tersebut, ketaksamaan Cauchy--Schwarz
berlaku. Terdapat banyak contoh lain yang memenuhi sifat-sifat ini selain vektor dalam
$\mathbb{R}^{n}$.

Ketaksamaan Cauchy--Schwarz memberikan bukti lain untuk \textbf{ketaksamaan
segitiga} bagi
\index{ketaksamaan segitiga}jarak dalam $\mathbb{R}^{n}$.

\begin{theorem}{Ketaksamaan Segitiga}\label{thm:triangleinequalitydotproduct}
Untuk $\vect{u},\vect{v}\in \mathbb{R}^{n}$,
\begin{equation}
\vectlength \vect{u}+\vect{v}\vectlength \leq \vectlength \vect{u}\vectlength +\vectlength \vect{v}
\vectlength  \label{triangleineq1}
\end{equation}
dan kesamaan berlaku jika dan hanya jika salah satu vektor merupakan kelipatan skalar tak negatif dari vektor lainnya.

Selain itu,
\begin{equation}
| \vectlength \vect{u}\vectlength -\vectlength \vect{v}\vectlength | \leq
\vectlength \vect{u}-\vect{v}\vectlength  \label{triangleineq2}
\end{equation}
\end{theorem}

\ifdefined\showproofs
\begin{proof} Berdasarkan sifat-sifat hasil kali titik dan ketaksamaan Cauchy--Schwarz,
\begin{eqnarray*}
\vectlength \vect{u}+\vect{v}\vectlength ^{2} &=& \left( \vect{u}+\vect{v}\right) \dotprod \left( \vect{u}+\vect{v}\right) \\
& =&\left( \vect{u}\dotprod \vect{u}\right) +\left( \vect{u}\dotprod \vect{v}\right) +\left(\vect{v}\dotprod \vect{u}\right) +\left( \vect{v}\dotprod \vect{v}\right) \\
&=&\vectlength \vect{u}\vectlength ^{2}+2\left( \vect{u}\dotprod \vect{v}\right)+\vectlength \vect{v}\vectlength ^{2} \\
&\leq& \vectlength \vect{u}\vectlength ^{2}+2\left\vert \vect{u}\dotprod \vect{v}\right\vert +\vectlength \vect{v}\vectlength ^{2} \\
&\leq& \vectlength \vect{u}\vectlength ^{2}+2\vectlength \vect{u}\vectlength \vectlength \vect{v}\vectlength +\vectlength \vect{v}\vectlength ^{2} =\left( \vectlength \vect{u}\vectlength +\vectlength \vect{v}\vectlength\right) ^{2}
\end{eqnarray*}
Jadi,
\begin{equation*}
\vectlength \vect{u}+\vect{v}\vectlength ^{2} \leq \left( \vectlength \vect{u}\vectlength +\vectlength \vect{v}\vectlength
\right) ^{2}
\end{equation*}
Dengan mengambil akar kuadrat kedua ruas, Anda memperoleh \ref{triangleineq1}.

Masih perlu dipertimbangkan kapan kesamaan terjadi. Misalkan $\vect{u} = \vect{0}$.
Maka $\vect{u} = 0 \vect{v}$ dan klaim mengenai kapan
kesamaan terjadi telah terverifikasi. Argumen yang sama berlaku jika $\vect{v} = \vect{0}$.
Oleh karena itu, dapat diasumsikan bahwa kedua vektor
tidak nol. Untuk memperoleh kesamaan dalam \ref{triangleineq1} di atas, Teorema \ref{thm:cauchyschwarzinequality}
menyiratkan bahwa salah satu vektor harus merupakan kelipatan
vektor lainnya. Misalkan $\vect{v}= k \vect{u}$. Jika $k <0$, maka kesamaan
tidak dapat terjadi dalam \ref{triangleineq1} karena dalam kasus ini
\begin{equation*}
\vect{u}\dotprod \vect{v} =k \vectlength \vect{u}\vectlength
^{2}<0<\left| k \right| \vectlength \vect{u}\vectlength ^{2}=\left| \vect{u}\dotprod \vect{v}\right|
\end{equation*}
Oleh karena itu, $k \geq 0.$

Untuk memperoleh bentuk lain ketaksamaan segitiga, tuliskan
\begin{equation*}
\vect{u}=\vect{u}-\vect{v}+\vect{v}
\end{equation*}
sehingga
\begin{align*}
\vectlength \vect{u}\vectlength & =\vectlength \vect{u}-\vect{v}+\vect{v}\vectlength \\
& \leq \vectlength \vect{u}-\vect{v}\vectlength +\vectlength \vect{v}\vectlength
\end{align*}
Oleh karena itu,
\begin{equation}
\vectlength \vect{u}\vectlength -\vectlength \vect{v}\vectlength \leq \vectlength \vect{u}-\vect{v}
\vectlength  \label{triangleineq3}
\end{equation}
Demikian pula,
\begin{equation}
\vectlength \vect{v}\vectlength -\vectlength \vect{u}\vectlength \leq \vectlength \vect{v}-\vect{u}
\vectlength =\vectlength \vect{u}-\vect{v}\vectlength  \label{triangleineq4}
\end{equation}
Dari \ref{triangleineq3} dan \ref{triangleineq4}, diperoleh bahwa \ref{triangleineq2} berlaku. Hal ini
karena $\left| \vectlength \vect{u}\vectlength -\vectlength \vect{v}\vectlength
\right| $ sama dengan ruas kiri dari \ref{triangleineq3} atau \ref{triangleineq4}, dan
dalam kedua keadaan, $\left| \vectlength \vect{u}\vectlength -\vectlength \vect{v}\vectlength
\right| \leq \vectlength \vect{u}-\vect{v}\vectlength $.
\end{proof}
\fi
